\par\Needspace{6\baselineskip}\medskip
\begin{center}\hypertarget{native-198-article148}{}\pdfbookmark[2]{148.}{article-148}148.\end{center}
Whenever the elements to be derived from three observations are still entirely unknown (the case to which our method is particularly suited), in the first hypothesis we take for $P$ and $Q$ the approximate values $\dfrac{\theta''}{\theta}$ and $\theta\theta''$, as already noted. Here $\theta$ and $\theta''$ are provisionally derived from the uncorrected time intervals. Expressing the ratios of these intervals to the corrected ones by $\mu:1$ and $\mu'':1$ respectively, we have in the first hypothesis
\[
\begin{aligned}
X&=\log\mu-\log\mu''+\log\eta-\log\eta''\\
Y&=-\log\mu-\log\mu''-\log\eta-\log\eta''\\
&\quad+\mathrm{Comp.}\log\cos f+\mathrm{Comp.}\log\cos f'\\
&\quad+\mathrm{Comp.}\log\cos f''+2\log r'-\log r-\log r''.
\end{aligned}
\]
The logarithms of the quantities $\mu,\mu''$ are insignificant compared with the remaining terms. The two positive quantities $\log\eta$ and $\log\eta''$ cancel one another to some extent in $X$, especially when the time intervals are nearly equal. Thus $X$ takes a small value, sometimes positive and sometimes negative. In $Y$, on the other hand, the negative terms $\log\eta$ and $\log\eta''$ do offset the positive terms $\mathrm{Comp.}\log\cos f$, $\mathrm{Comp.}\log\cos f'$, $\mathrm{Comp.}\log\cos f''$ to some extent, though less fully, for the latter usually exceed the former appreciably. In general nothing can be determined about the sign of $\log\dfrac{r'r'}{rr''}$.

Whenever the heliocentric motion between the observations is moderate, it will seldom be necessary to proceed to a fourth hypothesis. Usually the third and often
\SourceLeaf{199}even the second will give sufficient precision, and sometimes we can be content with the numbers resulting from the first hypothesis itself. It is always helpful to consider how precise the observations are. It would be a thankless task to aim for a precision in the calculation a hundred or a thousand times greater than the observations allow. In such matters, frequent practice sharpens judgment better than rules, and experienced calculators will readily acquire a sense of when to stop.
